\section{Distribution in congruence classes}\label{cg:section}

For an odd squarefree integer $d$ define
\begin{equation}\label{cg:remainder}
 r(d)=\sum_{\substack{u\ge1\\2u+1\equiv0\pmod d}}
                    A(u)\Psi(u/x)-\frac{X_A}{\varphi(d)}.
\end{equation}
The following estimate supplies the distribution needed by the sieve.
Its proof uses only the prime slot in $Q_0$, the weight bounds already
proved, and the classical Siegel--Walfisz and multiplicative large
sieve estimates in
\cite[Section ``Notation and preliminary estimates'']{SmoothShifted}.

\begin{theorem}[Congruence distribution]\label{cg:distribution}
For every fixed $0<\vartheta<1/2$ and $D>0$,
\[
 \sum_{\substack{d\le x^\vartheta\\d\ {\rm odd\ and\ squarefree}}}
       |r(d)|\ll_{\vartheta,D}xL^{-D}+X_A L^{-18}.
\]
The constants may depend on the fixed weight geometry and $\Psi$,
but not on any later proxy precision. No effective bound on the
threshold for $x$ is asserted.
\end{theorem}

\begin{proof}
Write $Q=x^\vartheta$ and $w(u)=A(u)\Psi(u/x)$. The congruence in
Equation~\eqref{cg:remainder} is the single unit class
$a_d\equiv-2^{-1}\pmod d$. Character orthogonality gives the exact
identity
\begin{equation}\label{cg:characters}
 r(d)=\frac1{\varphi(d)}
       \sum_{\substack{\chi\bmod d\\\chi\ne\chi_0}}
          \overline{\chi(a_d)}\sum_u w(u)\chi(u)
       -\frac1{\varphi(d)}\sum_{(u,d)>1}w(u).
\end{equation}
We first bound the correction in the second term. Mertens' estimate
implies
\begin{equation}\label{cg:totient-sum}
 \sum_{h\le Q}\frac{\mu^2(h)}{\varphi(h)}
 \le\prod_{p\le Q}\left(1+\frac1{p-1}\right)\ll L.
\end{equation}
For a fixed $u$ in the effective support of $w$, we have
$x\le u\le2x$ and every prime divisor of $u$ is at least $L^{20}$.
Writing a squarefree $d$ divisible by $p$ as $ph$, and then using a
union bound, gives
\[
 \sum_{\substack{d\le Q\\d\ {\rm squarefree}\\(d,u)>1}}
       \frac1{\varphi(d)}
 \le\sum_{p\mid u}\frac1{p-1}
        \sum_{h\le Q}\frac{\mu^2(h)}{\varphi(h)}
 \ll L\sum_{p\mid u}\frac1{p-1}\ll L^{-18}.
\]
For the last inequality, $u$ has at most
$\log(2x)/(20T)$ distinct prime divisors. Since $w\ge0$, the total
principal correction is therefore $O(X_A L^{-18})$.

To treat the nonprincipal characters, delete one designated ordered
slot of $Q_0$ from the definition of $a$. Retain the original divisor
$(r_0!)^{j_x+1}$ in the resulting coefficient, denoted by $a^-(v)$.
If the residual exponents are $e_p$, then
\[
 a^-(v)=\frac1{r_0}\prod_{j=0}^{j_x}\prod_{p\in Q_j}\frac1{e_p!}
 \le\frac1{r_0},
\]
with $r_0-1$ slots in $Q_0$ and $r_0$ slots in every other band, and
$a^-(v)=0$ off this support. Partitioning ordered tuples by the
deleted prime gives, including when primes repeat,
\[
 a(v)=\sum_{\substack{p\in Q_0\\p\mid v}}a^-(v/p).
\]
The $Q_0$-primes lie outside $\mathcal P$, so removing that slot
changes neither $W_1$ nor $\mathcal E$. Consequently
\begin{equation}\label{cg:prime-convolution}
 A(u)=\sum_{\substack{pr=u\\p\in Q_0}}b(r),
 \qquad
 b(r)=W_1(r)\frac{1-\mathcal E(r)}2a^-(r_*),
 \qquad |b(r)|\le L^{B_0},
\end{equation}
where $B_0$ is fixed by the weight geometry.

Split $p,r$ into dyads of sizes $P,R$. Only $O(L)$ pairs of dyads
contribute, and on these $PR\asymp x$. Since
$x^{c_1s'}\le p\le x^{c_2s'}$, there is a fixed $c>0$, depending
only on Equation~\eqref{wt:geometry}, such that
\begin{equation}\label{cg:factor-ranges}
 P,R\ge x^c
\end{equation}
for every contributing pair and all sufficiently large $x$.
For example any $c<\min(c_1s',1-c_2s')$ suffices.

Every character modulo an odd squarefree $d$ is induced by a unique
primitive character $\chi_f$ of conductor $f\mid d$. Write $d=fh$;
then $(f,h)=1$ and $\varphi(d)=\varphi(f)\varphi(h)$. For such an
induced character,
\[
 \chi(pr)=\chi_f(p)\chi_f(r)
                  \1_{\{(p,h)=1\}}\1_{\{(r,h)=1\}}.
\]
The nonprincipal characters have $f>1$; there is no multiplicity in
this parametrization. After applying the triangle inequality in
Equation~\eqref{cg:characters}, it remains to bound sums of the form
\begin{equation}\label{cg:primitive-sum}
 \begin{split}
 \sum_{\substack{h\le Q\\h\ {\rm odd\ and\ squarefree}}}
 \frac1{\varphi(h)}
 \sum_{\substack{1<f\le Q/h\\f\ {\rm odd\ and\ squarefree}\\(f,h)=1}}
 &\frac1{\varphi(f)}\sum_{\chi_f\bmod f}^{*}\\
 &\quad\cdot
 \left|\sum_{\substack{p\in Q_0,\ p\asymp P\\r\asymp R}}
       b(r)\chi_f(pr)\1_{\{(pr,h)=1\}}\Psi(pr/x)\right|.
 \end{split}
\end{equation}
Here and below the asterisk denotes primitive characters. Modulus
restrictions may be discarded in nonnegative majorants, which will
be useful when applying the large sieve.

Choose a fixed $K$ later, and first consider the conductors
$f>L^K$. Put $\psi(v)=\Psi(e^v)$ and
$\widehat\psi(t)=\int_{\mathbb R}\psi(v)e^{-itv}\,dv$.
Fourier inversion gives
\[
 \Psi(pr/x)=\frac1{2\pi}\int_{\mathbb R}
             \widehat\psi(t)x^{-it}p^{it}r^{it}\,dt,
 \qquad \int_{\mathbb R}|\widehat\psi(t)|\,dt\ll_\Psi1.
\]
There is no frequency truncation error. For fixed $h,t$, define
\[
 \mathcal P_h(\chi,t)=
    \sum_{\substack{p\in Q_0,\ p\asymp P\\(p,h)=1}}\chi(p)p^{it},
 \qquad
 \mathcal R_h(\chi,t)=
    \sum_{\substack{r\asymp R\\(r,h)=1}}b(r)\chi(r)r^{it}.
\]
Their coefficient squared norms are at most $O(P)$ and
$O(RL^{2B_0})$, respectively, uniformly in $h,t$. On a conductor
dyad $F<f\le2F$, Cauchy--Schwarz followed by the primitive
multiplicative large sieve therefore yields
\begin{align}
 \sum_{F<f\le2F}\frac1{\varphi(f)}
       \sum_{\chi_f\bmod f}^{*}
       |\mathcal P_h(\chi_f,t)\mathcal R_h(\chi_f,t)|
 &\ll \frac{L^{B_0}}{F}
       \sqrt{(P+F^2)(R+F^2)PR} \notag\\
 &\ll xL^{B_0}
       \left(\frac1F+\frac1{\sqrt P}+\frac1{\sqrt R}
                        +\frac F{\sqrt x}\right).
       \label{cg:large-sieve-bound}
\end{align}
Indeed $1/\varphi(f)\ll F^{-1}f/\varphi(f)$ on this dyad, giving
exactly the large sieve's primitive-character weight. The masks
$\1_{(p,h)=1},\1_{(r,h)=1}$ do not enlarge the coefficient norms.

Integrate against $|\widehat\psi|$, sum the conductor and factor
dyads, and use Equation~\eqref{cg:totient-sum} for the remaining
$h$-sum. Since $F\le Q$, Equations~\eqref{cg:factor-ranges} and
\eqref{cg:large-sieve-bound} give a total at most
\begin{equation}\label{cg:large-conductors}
 xL^{B_2}
       \left(L^{-K}+x^{-c/2}+x^{\vartheta-1/2}\right)
\end{equation}
for a fixed $B_2$ independent of $K$. One can use dyads intersected
with $f>L^K$, so their lower endpoints are at least $L^K/2$.
In particular the summation over $h$ costs a logarithmic factor,
not a factor $Q$. Taking $K>B_2+D+2$ makes
Equation~\eqref{cg:large-conductors} $O(xL^{-D})$, because
$\vartheta<1/2$.

It remains to consider $1<f\le2L^K$. For fixed $r\asymp R$, apply
Siegel--Walfisz to the nonprincipal character $\chi_f$ on the
prime interval $Q_0\cap[P,2P)$. Since $P\ge x^c$, the modulus is
bounded by a fixed power of $\log P$. Summing the residue-class
estimate against $\chi_f$ and asking for a stronger initial saving
absorbs the at most $\varphi(f)$ classes. Partial summation, using
the uniformly bounded supremum and total variation of
$p\mapsto\Psi(pr/x)$ on this interval, gives, for any fixed $M$,
\begin{equation}\label{cg:small-prime-sum}
 \sum_{\substack{p\in Q_0\\p\asymp P}}
       \chi_f(p)\Psi(pr/x)\ll_{K,M}P L^{-M}.
\end{equation}
This use of Siegel--Walfisz includes possible exceptional real
characters; its potentially ineffective constants cause no problem.

Imposing $(p,h)=1$ deletes only $O(1)$ terms in
Equation~\eqref{cg:small-prime-sum}: each deleted prime is at least
$x^c$, while $h\le x^\vartheta$, so there are at most
$\vartheta/c+O(1)$ such primes. Thus the inner double sum in
Equation~\eqref{cg:primitive-sum} is at most
\[
 RL^{B_0}\bigl(PL^{-M}+O(1)\bigr)
 \ll xL^{B_0-M}+x^{1-c}L^{B_0}.
\]
The other mask, $(r,h)=1$, only decreases this absolute bound.
For the small conductors,
\[
 \sum_{1<f\le2L^K}\frac1{\varphi(f)}
                         \sum_{\chi_f\bmod f}^{*}1\ll L^K.
\]
Combining this with Equation~\eqref{cg:totient-sum} and the number of
factor dyads bounds the small-conductor contribution by
\begin{equation}\label{cg:small-conductors}
 xL^{B_3+K-M}+x^{1-c}L^{B_3+K}
\end{equation}
for a fixed $B_3$ independent of $K,M$. Choose
$M>B_3+K+D+2$. Equation~\eqref{cg:small-conductors} is then
$O(xL^{-D})$. Together with the large-conductor estimate and the
principal correction, this proves the theorem.
\end{proof}
